\scanpage{547}
\mspage{458}

\uline{Lemme} 4.1\deleted{\unsure{6}}5.8.2. Soit $(E_i)_{i\in\unsure{J|I}}$ un système inductif d'ensembles, indexé par un ens. ordonné \inserted{$I$} grand devant $\pi_0$. On suppose qu'il existe $i_0\in I$ tel que

a) $i\geq i_0$ implique $\operatorname{card} E_i=\pi_0$, et

b) $j\geq i\geq i_0$ implique que $\deleted{\unsure{\alpha}}\unsure{u|\alpha}_{ji}:E_i\mathbin{\unsure{-|\to}} E_j$ est un mono. Supposons de plus

c) $\forall\,\deleted{\unsure{j_0|j}}i\deleted{\unsure{\in|>}}\geq \unsure{I_0|i_0}$, $\exists\, j>\deleted{\unsure{j_0|j}}i$ tel que $\deleted{\unsure{\alpha}}u_{j\deleted{\unsure{j_0|i_0}}i}$ ne soit \uline{pas} iso. (i.e. ne soit pas surjectif).

\deleted{Alors} Soit $E_\infty=\underset{i\in I}{\unsure{\underline{\lim}|\varinjlim}}\,E_i$. Alors
\[
\operatorname{card} E_\infty>\pi_0.
\]

Ce lemme implique donc l'\uline{équivalence} de 4.8.13 et du ``lemme-clef'' 4.15.8.1.

Pour prouver 4.15.8.2, nous allons prouver quelque chose d'un peu plus précis :

\uline{Corollaire} 4.15.8.3. Sous les conditions de \unsure{4.15.8.2|4.35.8.2}, soit $\alpha_0$ le premier ordinal tel que $\operatorname{card}\alpha_0>\pi_0$, et soit $I_{\alpha_0}$ l'ens des ordinaux $\beta<\alpha$ (ensemble \inserted{ordonné} qui est grand devant $\pi_0$). Il existe alors une application

\inserted{\unsure{str.}} croissante $I_{\alpha_0}\longrightarrow I$, $\alpha\longmapsto i_\alpha$, \inserted{appliquant $0$ dans $i_0$}
% archive sheet number: 546
